% !TeX program = xelatex
\documentclass[a4paper,oneside,11pt]{ctexrep}

\usepackage[a4paper,left=0.78in,right=0.78in,top=2cm,bottom=2cm]{geometry}
\usepackage{fontspec}
\setmainfont{TeX Gyre Heros}
\setCJKmainfont{Microsoft YaHei}
\usepackage{fancyhdr}
\usepackage[linesnumbered,ruled,vlined]{algorithm2e}
\usepackage{amsmath,amssymb,mathtools,bm}
\usepackage{siunitx}
\usepackage{booktabs,longtable,multirow,makecell,array,tabularx}
\usepackage{graphicx,float,adjustbox}
\usepackage{xcolor,colortbl}
\usepackage{listings}
\usepackage{tikz}
\usetikzlibrary{arrows.meta,positioning,shapes.geometric}
\usepackage[font=small,labelfont=bf]{caption}
\usepackage[colorlinks=true,linkcolor=purple,urlcolor=blue,citecolor=red]{hyperref}
\usepackage{url}

\definecolor{HeaderBlue}{HTML}{17365D}
\definecolor{LightBlue}{HTML}{DCE6F1}
\definecolor{LightOrange}{HTML}{FCE5CD}
\definecolor{LightGreen}{HTML}{D9EAD3}
\definecolor{LightGray}{HTML}{F3F3F3}

\renewcommand{\arraystretch}{1.25}
\setlength{\tabcolsep}{7pt}
\setlength{\parindent}{2em}
\setlength{\parskip}{0.25em}
\captionsetup[figure]{name={图}}
\captionsetup[table]{name={表}}

\pagestyle{fancy}
\fancyhf{}
\fancyfoot[C]{\small 基带电路自动评估测试框架\quad\thepage}
\renewcommand{\headrulewidth}{0pt}
\renewcommand{\footrulewidth}{0pt}

\lstdefinelanguage{json}{
    basicstyle=\ttfamily\small,
    numbers=left,
    numberstyle=\tiny\color{gray},
    stepnumber=1,
    numbersep=8pt,
    showstringspaces=false,
    breaklines=true,
    frame=single,
    rulecolor=\color{gray!45},
    backgroundcolor=\color{LightGray},
    literate=
     *{0}{{{\color{blue!65!black}0}}}{1}
      {1}{{{\color{blue!65!black}1}}}{1}
      {2}{{{\color{blue!65!black}2}}}{1}
      {3}{{{\color{blue!65!black}3}}}{1}
      {4}{{{\color{blue!65!black}4}}}{1}
      {5}{{{\color{blue!65!black}5}}}{1}
      {6}{{{\color{blue!65!black}6}}}{1}
      {7}{{{\color{blue!65!black}7}}}{1}
      {8}{{{\color{blue!65!black}8}}}{1}
      {9}{{{\color{blue!65!black}9}}}{1}
}

\lstdefinestyle{command}{
    basicstyle=\ttfamily\small,
    breaklines=true,
    frame=single,
    rulecolor=\color{gray!45},
    backgroundcolor=\color{LightGray},
    showstringspaces=false
}

\newcommand{\metricheader}[1]{\textbf{#1}}

\title{基带电路自动评估测试框架}
\author{}
\date{2026年9月}

\begin{document}

\begin{titlepage}
    \centering
    \vspace*{3.0cm}
    {\Huge\bfseries 基带电路自动评估测试框架\par}
    \vspace{1.2cm}
    {\Large 第一版\par}
    \vspace{2.2cm}
    \begin{tabular}{ll}
        \toprule
        文档类型 & 测试框架与接入说明 \\
        覆盖范围 & Metrics框架及已有五个模块 \\
        当前版本 & 0.1 \\
        日期 & 2026年9月 \\
        \bottomrule
    \end{tabular}
    \vfill
\end{titlepage}

\tableofcontents
\thispagestyle{empty}
\clearpage
\setcounter{page}{1}

\chapter{测试目标和模块范围}

\section{目的}

本文档介绍基带电路自动评估框架的使用流程、软件接口及结果输出方式。用户首先通过JSON配置文件设置待测模块的结构参数、量化参数和时钟参数；框架据此生成面积、延迟、吞吐率和硬件计算复杂度的预测结果。参考值由传统评估流程获得，其中Design Compiler用于生成面积及综合时序结果，Icarus Verilog用于执行RTL仿真，并结合C/C++黄金模型验证功能正确性以及获取周期级延迟和输出间隔。框架最终统一输出预测值、参考值、预测误差、评估时间和速度提升倍数。

\section{评估与验收指标}

框架通过Predict、Validate和Evaluate三条处理链完成自动评估。Predict链根据模块配置和时钟参数生成各项指标的预测值；Validate链调用C/C++参考程序、RTL仿真工具和逻辑综合工具，完成功能核验并获得对应参考值；Evaluate链汇总两条链的结果，计算预测误差、评估时间和速度提升倍数。

功能正确性是开展指标评估的前提。只有当RTL输出与C/C++黄金输出逐项一致时，对应Case才进入误差、时间和速度提升倍数的统计；若功能验证未通过，则记录验证失败原因，但不将该Case纳入指标验收结果。

在功能验证通过的基础上，框架按照表中规定的计算口径和验收要求对自动评估结果进行判定。

\section{评估指标}

\begin{longtable}{p{0.34\textwidth}p{0.20\textwidth}p{0.36\textwidth}}
    \toprule
    \metricheader{评估指标} & \metricheader{验收要求} & \metricheader{判定数据} \\
    \midrule
    \endfirsthead
    \toprule
    \metricheader{功能指标} & \metricheader{验收要求} & \metricheader{判定数据} \\
    \midrule
    \endhead
    基带电路自动评估时间 & 不超过\SI{10}{\second} & 四项指标预测时间之和 \\
    基带电路自动评估速度提升 & 不低于10倍 & 参考评估时间与自动评估时间之比 \\
    基带电路自动评估误差 & 不超过25\% & 预测值与参考值的绝对相对误差 \\
    定点计算复杂度评估误差 & 不超过15\% & 预测与参考$\mathrm{GE\cdot cycle}$的绝对相对误差 \\
    \bottomrule
\end{longtable}

\section{模块范围}

框架计划覆盖FxMatch、Delay、Neg、AbsMux（2:1和N:M）、Add、Sub、Mul、Comp、Mac Systolic Array、CompTree、AdderTree、BCB、Perm、Counter、Early Stop、Hard Decision、5G PUSCH CE、INSA-MMSE MIMO Detector和BP Polar Decoder。

本阶段先展开框架中已有的Add、Mul、5G PUSCH CE、INSA-MMSE MIMO Detector和BP Polar Decoder。每个模块选取一个已有Case说明预测和验证过程，其余模块后续按照相同接口接入。

\section{主要符号}

\begin{longtable}{lp{0.72\textwidth}}
    \toprule
    \metricheader{符号} & \metricheader{含义} \\
    \midrule
    \endfirsthead
    \toprule
    \metricheader{符号} & \metricheader{含义} \\
    \midrule
    \endhead
    $A_{\mathrm{pred}}$ & 面积预测值，单位为$\mu\mathrm{m}^{2}$ \\
    $A_{\mathrm{ref}}$ & 综合报告或参考数据库给出的面积 \\
    $L_{\mathrm{pred}}$ & 延迟预测值，单位为cycle \\
    $L_{\mathrm{RTL}}$ & 从RTL波形测得的延迟，单位为cycle \\
    $I_{\mathrm{pred}}$ & 预测输出间隔，单位为cycle \\
    $I_{\mathrm{RTL}}$ & RTL实测输出间隔，单位为cycle \\
    $B_{\mathrm{valid}}$ & 一次有效输出包含的有效比特数 \\
    $A_{\mathrm{GE}}$ & 一个基准等效门的面积 \\
    \bottomrule
\end{longtable}

\chapter{Metrics框架和模块接入}

前述预测、验证和指标计算流程由Metrics框架统一组织。该框架通过标准化配置文件接收测试参数，以模块Adapter屏蔽不同算法模块在运行环境、调用命令和结果格式方面的差异，并将各模块内部的预测程序、参考程序、RTL仿真流程和综合流程接入统一的Predict、Validate和Evaluate接口。本章首先介绍Metrics框架的软件结构和调用关系，随后说明各类基带模块的接入方式及结果返回协议。
\section{现有软件结构}

仓库根目录中的\texttt{metrics\_framework}负责统一命令行入口、模块注册、配置解析、子进程调度、协议检查和结果渲染。各算法模块保留各自的源码目录、Python运行环境、RTL生成程序和参考验证工具，并通过Adapter接入Metrics框架，从而在不改变模块内部实现的情况下完成统一调用和结果汇总。

\begin{longtable}{p{0.22\textwidth}p{0.31\textwidth}p{0.37\textwidth}}
    \toprule
    \metricheader{组成} & \metricheader{仓库位置} & \metricheader{职责} \\
    \midrule
    \endfirsthead
    \toprule
    \metricheader{组成} & \metricheader{仓库位置} & \metricheader{职责} \\
    \midrule
    \endhead
    CLI入口 & \path{metrics_framework/cli.py} & 解析模块名、模式和Case参数，输出机器可解析JSON \\
    Core调度层 & \path{metrics_framework/core.py} & 解析配置、启动Adapter、校验协议、计算误差和写入结果 \\
    模块注册表 & \path{metrics_framework/registry.json} & 记录模块目录、Adapter、解释器、配置规则和能力 \\
    公共协议封装 & \path{metrics_framework/adapters/common.py} & 统一Adapter的参数读取、返回码和协议JSON \\
    模块Adapter & \path{metrics_framework/adapters/*.py} & 将统一动作转换为模块内部预测和验证调用 \\
    模块实现 & 各模块独立目录 & 保存模型、RTL生成器、行为参考程序和仿真工作区 \\
    \bottomrule
\end{longtable}

\section{当前注册模块}

\begin{longtable}{llll}
    \toprule
    \metricheader{模块名} & \metricheader{模块目录} & \metricheader{Adapter} & \metricheader{默认Case} \\
    \midrule
    \endfirsthead
    \toprule
    \metricheader{模块名} & \metricheader{模块目录} & \metricheader{Adapter} & \metricheader{默认Case} \\
    \midrule
    \endhead
    \texttt{add} & \path{Generator/BasicTests/Add} & \path{adapters/add.py} & 1--5 \\
    \texttt{mul} & \path{Generator/BasicTests/Mul} & \path{adapters/mul.py} & 1--5 \\
    \texttt{ls} & \path{Generator/LSCE} & \path{adapters/ls.py} & 1--5 \\
    \texttt{mimo} & \path{Generator/MIMODetector} & \path{adapters/mimo.py} & 1--5 \\
    \texttt{bp} & \path{BPPredIter/BP_Evaluation} & \path{adapters/bp.py} & 1--5 \\
    \bottomrule
\end{longtable}

注册表不要求不同模块使用同一套配置字段。它只统一模块定位、配置文件查找、输出目录和Adapter调用方式。新模块接入时实现自己的Predict和Validate函数，并在\texttt{registry.json}增加一条模块记录。

\section{对外调用方法}

命令行调用格式为：

\begin{lstlisting}[style=command]
python -m metrics_framework <module> <predict|evaluate> [case-or-config]
\end{lstlisting}

例如：

\begin{lstlisting}[style=command]
python -m metrics_framework add predict 1
python -m metrics_framework add evaluate 1
python -m metrics_framework mimo evaluate 1
python -m metrics_framework bp evaluate BPPredIter/BP_Evaluation/config/config1.json
\end{lstlisting}

省略Case参数时，框架按照注册表中的\texttt{default\_cases}运行默认Case。数字参数按模块的\texttt{config\_pattern}解析为配置文件路径；路径参数则直接指向指定的\texttt{config.json}。

Python API可直接调用：

\begin{lstlisting}[language=Python,style=command]
from metrics_framework import predict, evaluate

prediction = predict("add", "1")
evaluation = evaluate("mimo", "1")
\end{lstlisting}

\section{内部接入关系}

统一入口取得模块名后，从\texttt{registry.json}读取模块根目录、Adapter路径、配置规则和Python解释器环境变量。Core调度层随后在模块自己的工作目录中启动Adapter子进程，避免不同模块的依赖、同名包或ABI相互干扰。

Adapter内部提供两个动作：

\begin{itemize}
    \item \texttt{predict CONFIG}：调用模块预测函数，返回四项预测指标及分项计时；
    \item \texttt{validate CONFIG}：调用行为参考、RTL仿真和综合参考，返回真实值及验证证据。
\end{itemize}

\texttt{validate}是框架内部动作，不是统一CLI面向用户公开的运行模式。用户执行\texttt{evaluate}时，Core先调用Adapter的\texttt{predict}，再调用同一Adapter的\texttt{validate}，最后组合两次返回值。

\section{Predict Validate和Evaluate}

三者的关系为：

\begin{equation}
    \mathrm{Evaluate}
    =\mathrm{Predict}+\mathrm{Validate}+\mathrm{ResultCalculation}.
\end{equation}

Predict不运行RTL仿真和逻辑综合，也不读取待比较的真实值。Validate负责功能验证、RTL时序测量和综合参考提取。Evaluate负责顺序调用、配置一致性校验、误差计算、时间统计和结果保存。

预测和验证协议均包含模块名、动作、配置路径、配置摘要、证据路径、指标来源和未舍入指标值。Core比较两次调用的配置摘要，防止预测和验证使用不同配置。

若Validate返回参考数据不可用，Evaluate不输出不完整的\texttt{evaluation.json}，并在诊断目录保存预测协议结果和失败原因。批量Evaluate中任一Case失败时，批次不输出部分JSON。

\section{框架调用流程}

\begin{figure}[H]
    \centering
    \resizebox{0.96\textwidth}{!}{%
    \begin{tikzpicture}[
        node distance=8mm and 14mm,
        box/.style={draw=HeaderBlue,rounded corners,thick,align=center,minimum height=10mm,text width=44mm},
        predict/.style={box,fill=LightBlue},
        validate/.style={box,fill=LightOrange},
        result/.style={box,fill=LightGreen},
        line/.style={-Latex,thick,draw=gray!75}
    ]
        \node[predict] (config) {模块\texttt{config.json}\\与时钟参数};
        \node[result,below=of config] (eval) {Evaluate统一入口};
        \node[predict,below left=13mm and 17mm of eval] (pred) {Predict链\\面积 延迟 吞吐率 复杂度};
        \node[validate,below right=13mm and 17mm of eval] (val) {Validate链\\C/C++ RTL 综合结果};
        \node[predict,below=of pred] (predout) {预测结果与分项计时\\$A_{pred},L_{pred},R_{pred},C_{pred}$};
        \node[validate,below=of val] (valout) {功能匹配与参考结果\\$A_{ref},L_{RTL},I_{RTL}$\\$R_{ref},C_{ref}$};
        \node[result,below=19mm of eval,yshift=-42mm] (calc) {误差与时间计算\\自动评估时间 速度提升};
        \node[result,below=of calc] (output) {\texttt{evaluation.json}与汇总表};
        \draw[line] (config) -- (eval);
        \draw[line] (eval) -- (pred);
        \draw[line] (eval) -- (val);
        \draw[line] (pred) -- (predout);
        \draw[line] (val) -- (valout);
        \draw[line] (predout) -- (calc);
        \draw[line] (valout) -- (calc);
        \draw[line] (calc) -- (output);
    \end{tikzpicture}%
    }
    \caption{Metrics框架调用与模块接入关系}
    \label{fig:metrics-flow}
\end{figure}

\section{输出文件}

\begin{longtable}{p{0.24\textwidth}p{0.48\textwidth}p{0.19\textwidth}}
    \toprule
    \metricheader{结果} & \metricheader{内容} & \metricheader{生成条件} \\
    \midrule
    \endfirsthead
    \toprule
    \metricheader{结果} & \metricheader{内容} & \metricheader{生成条件} \\
    \midrule
    \endhead
    \texttt{prediction.json} & 四项预测值、分项预测时间和自动评估总时间 & Predict成功 \\
    Validate协议结果 & 功能比对、RTL延迟、输出间隔、综合面积和参考时间 & Validate成功 \\
    \texttt{evaluation.json} & 预测值、参考值、误差、时间和速度提升倍数 & 两条链均成功 \\
    诊断记录 & 配置、协议、功能比对、仿真或参考数据错误 & Evaluate未完成 \\
    \bottomrule
\end{longtable}

\chapter{测试输入 Validate验证和指标计算}

\section{测试输入}

\subsection{模块配置文件}

各模块使用自己的\texttt{config.json}描述设计参数。以Add模块Case 1为例：

\begin{lstlisting}[language=json]
{
  "input_1": {"bitwidth": 1, "fractional_width": -1, "signed": false},
  "input_2": {"bitwidth": 1, "fractional_width": -1, "signed": false},
  "output": {"bitwidth": 1, "fractional_width": 0, "signed": false},
  "n_pipeline": 1,
  "if_rst_n": false,
  "clock": {"period_ns": 10.0}
}
\end{lstlisting}

该配置给出两路输入和输出的总位宽、小数位宽及符号属性，同时给出流水级数、复位形式和时钟周期。系统级模块还包含天线数量、并行度、量化方式、调制阶数、码长、码率和迭代次数等参数。

\subsection{时钟参数}

若测试条件给出时钟频率，框架先换算为时钟周期：

\begin{equation}
    T_{\mathrm{clk}}(\mathrm{ns})
    =\frac{1000}{f_{\mathrm{clk}}(\mathrm{MHz})}.
\end{equation}

延迟以cycle输出，吞吐率使用时钟周期换算为单位时间内的有效输出量。

\section{Validate功能验证和参考值提取}

Validate使用同一份配置和输入数据分别运行C/C++行为模型与RTL仿真程序。输入、黄金输出和RTL输出通过帧编号、样本编号或矩阵元素位置建立一一对应关系。

\begin{algorithm}[H]
    \caption{Validate通用功能验证流程}
    \KwIn{模块名称$Module$，模块配置$Config$}
    \KwOut{功能匹配结论及参考指标}
    $ConfigData \leftarrow LoadAndCheck(Config)$\;
    $InputData \leftarrow GenerateInput(ConfigData)$\;
    $CPPOutput \leftarrow RunCPPReference(ConfigData, InputData)$\;
    $RTLFiles \leftarrow GenerateRTL(Module, ConfigData)$\;
    $(RTLOutput, Waveform, T_{RTL}) \leftarrow RunRTLSimulation(RTLFiles, InputData)$\;
    \If{$Count(RTLOutput) \neq Count(CPPOutput)$}{
        \Return Fail\;
    }
    \ForEach{$SampleID$}{
        $CPPValue \leftarrow NormalizeFixedPoint(CPPOutput[SampleID])$\;
        $RTLValue \leftarrow NormalizeFixedPoint(RTLOutput[SampleID])$\;
        \If{$RTLValue \neq CPPValue$}{
            $RecordMismatch(SampleID, CPPValue, RTLValue)$\;
        }
    }
    \If{$MismatchCount > 0$}{
        \Return Fail\;
    }
    $L_{RTL} \leftarrow MeasureLatency(Waveform)$\;
    $I_{RTL} \leftarrow MeasureOutputInterval(Waveform)$\;
    $(A_{ref},T_{syn}) \leftarrow ReadSynthesisResult(ConfigData)$\;
    $R_{ref} \leftarrow CalculateThroughput(I_{RTL},ConfigData)$\;
    $C_{ref} \leftarrow CalculateComplexity(A_{ref},L_{RTL})$\;
    \Return Pass及全部参考值\;
\end{algorithm}

定点结果比较前，需要统一符号、补码、总位宽、小数位宽、舍入、截断、溢出、饱和及输出顺序。只有全部RTL输出与C/C++黄金输出一致，波形测得的延迟和输出间隔才作为有效参考值。

\section{面积}

Predict调用面积模型，根据位宽、流水级数、并行度和结构规模得到预测面积。Validate从逻辑综合报告或已有综合数据库取得参考面积和综合时间。面积单位为$\mu\mathrm{m}^{2}$。

\begin{align}
    A_{\mathrm{pred}} &= F_{\mathrm{area}}(Config),\\
    E_A &= \frac{|A_{\mathrm{pred}}-A_{\mathrm{ref}}|}{A_{\mathrm{ref}}}\times100\%.
\end{align}

\section{延迟}

预测延迟由流水级数、结构调度或迭代次数模型得到。RTL延迟从输入被接收的时钟边沿开始，到对应输出有效的时钟边沿结束。预测值和参考值均以cycle表示。

\begin{align}
    L_{\mathrm{RTL}} &= k_{\mathrm{out}}-k_{\mathrm{in}},\\
    E_L &= \frac{|L_{\mathrm{pred}}-L_{\mathrm{RTL}}|}{L_{\mathrm{RTL}}}\times100\%.
\end{align}

\section{吞吐率}

预测吞吐率使用预测输出间隔，参考吞吐率使用RTL波形实测输出间隔。后期各模块统一使用bit吞吐率，单位为Gbps。当前程序中的Gframes/s和kChannels/s可根据一次有效输出包含的比特数换算。

\begin{align}
    R_{\mathrm{pred}} &= \frac{B_{\mathrm{valid}}}{T_{\mathrm{clk}}I_{\mathrm{pred}}},\\
    R_{\mathrm{RTL}} &= \frac{B_{\mathrm{valid}}}{T_{\mathrm{clk}}I_{\mathrm{RTL}}},\\
    E_T &= \frac{|R_{\mathrm{pred}}-R_{\mathrm{RTL}}|}{R_{\mathrm{RTL}}}\times100\%.
\end{align}

其中$B_{\mathrm{valid}}$是一次有效输出包含的有效比特数。$T_{\mathrm{clk}}$以ns表示时，吞吐率结果为Gbps。

\section{硬件计算复杂度}

硬件计算复杂度先将面积换算为等效门数量，再乘以完成一次计算所需的周期数：

\begin{align}
    C_{\mathrm{pred}} &= \frac{A_{\mathrm{pred}}}{A_{\mathrm{GE}}}L_{\mathrm{pred}},\\
    C_{\mathrm{ref}} &= \frac{A_{\mathrm{ref}}}{A_{\mathrm{GE}}}L_{\mathrm{RTL}},\\
    E_C &= \frac{|C_{\mathrm{pred}}-C_{\mathrm{ref}}|}{C_{\mathrm{ref}}}\times100\%.
\end{align}

复杂度单位为$\mathrm{GE\cdot cycle}$。该指标同时反映硬件资源和计算周期，适合分析串行、流水和并行结构的面积与延迟权衡。文献\cite{nguyen2025mlkem}使用面积时间积比较ML-KEM硬件架构效率；文献\cite{curlin2025sbox}从面积、关键路径延迟、功耗和安全等维度比较AES SBox实现。两项工作均说明硬件实现需要联合考察资源和时间。

$\mathrm{GE\cdot cycle}$不能替代功耗、能耗、最高工作频率或算法精度。不同设计进行横向比较时，应采用相同工艺、GE基准单元、功能规模和输入条件。

\section{时间和速度提升}

自动评估总时间为四项指标预测时间之和。参考时间由RTL验证时间和面积综合时间构成：

\begin{align}
    T_{\mathrm{auto}} &= T_A+T_L+T_T+T_C,\\
    T_{\mathrm{ref}} &= T_{\mathrm{syn}}+T_{\mathrm{RTL}},\\
    S_{\mathrm{total}} &= \frac{T_{\mathrm{ref}}}{T_{\mathrm{auto}}}.
\end{align}

\chapter{已有模块代表Case}

每个模块选择一个已有Case，按照配置、预测公式、Validate过程和结果记录展开。面积模型的内部系数由模块现有模型给出，本章使用函数形式表示。

\section{Add模块Case 1}

\subsection{配置}

两路输入均为1 bit无符号定点数，输入\texttt{fractional\_width}为$-1$；输出为1 bit无符号数；流水级数为1；时钟周期为\SI{10}{\nano\second}。

\subsection{预测公式}

\begin{align}
    A_{\mathrm{pred}} &= F_{\mathrm{ADD\_area}}(Config_{case1}),\\
    L_{\mathrm{pred}} &= n_{\mathrm{pipeline}}=1\ \mathrm{cycle},\\
    I_{\mathrm{pred}} &= 1\ \mathrm{cycle},\\
    R_{\mathrm{pred}} &= \frac{1\ \mathrm{bit}}{10\ \mathrm{ns}\times1}=0.1\ \mathrm{Gbps},\\
    C_{\mathrm{pred}} &= \frac{A_{\mathrm{pred}}}{A_{\mathrm{GE}}}\times1.
\end{align}

\subsection{Validate过程}

\begin{enumerate}
    \item 按Case 1配置生成连续多帧输入，并运行C/C++定点加法参考程序。
    \item 生成Add、FxMatch、Delay和testbench RTL，完成编译与仿真。
    \item 逐帧比较RTL输出和C/C++输出，从VCD测量延迟及相邻输出间隔。
    \item 读取对应综合面积和综合时间，计算参考吞吐率与$\mathrm{GE\cdot cycle}$。
    \item 记录四项误差、自动评估时间、参考时间和速度提升倍数。
\end{enumerate}

\section{Mul模块Case 1}

\subsection{配置}

两路输入均为4 bit有符号定点数，小数位宽为2；输出为7 bit有符号定点数，小数位宽为3；流水级数为1；时钟周期为\SI{5}{\nano\second}。

\subsection{预测公式}

\begin{align}
    A_{\mathrm{pred}} &= F_{\mathrm{MUL\_area}}(Config_{case1}),\\
    L_{\mathrm{pred}} &= 1\ \mathrm{cycle},\qquad I_{\mathrm{pred}}=1\ \mathrm{cycle},\\
    R_{\mathrm{pred}} &= \frac{7\ \mathrm{bits}}{5\ \mathrm{ns}\times1}=1.4\ \mathrm{Gbps},\\
    C_{\mathrm{pred}} &= \frac{A_{\mathrm{pred}}}{A_{\mathrm{GE}}}\times1.
\end{align}

\subsection{Validate过程}

Validate运行QuBLAS或C++乘法参考程序和Mul RTL仿真，逐帧比较乘法结果，检查符号、乘积位宽、小数位截断和溢出处理；随后从VCD取得实际延迟和输出间隔，从综合数据取得实际面积。

\section{5G PUSCH CE模块Case 1}

\subsection{配置}

发送和接收天线数均为16，发送并行度为4，接收并行度为16，乘法和Adder Tree流水级数分别为1和3，时钟周期为\SI{10}{\nano\second}。

\subsection{预测公式}

\begin{align}
    S_T &= \left\lceil\frac{16}{4}\right\rceil=4,\\
    L_{\mathrm{pred}} &= 1+3+4=8\ \mathrm{cycles},\\
    R_{\mathrm{pred}} &= \frac{10^6}{4\times10\times16}=1562.5\ \mathrm{kChannels/s},\\
    A_{\mathrm{pred}} &= F_{\mathrm{CE\_area}}(Config_{case1}),\\
    C_{\mathrm{pred}} &= \frac{A_{\mathrm{pred}}}{A_{\mathrm{GE}}}\times8.
\end{align}

本Case先保留程序当前的kChannels/s输出。确定每个信道估计结果的有效比特数后，再换算为Gbps。

\subsection{Validate过程}

Validate运行C/C++信道估计参考模型和RTL仿真，按照接收天线、发送天线及输出位置比较结果；从波形测量实际延迟和输出间隔，并从综合数据取得实际面积。

\section{INSA-MMSE MIMO Detector模块Case 1}

\subsection{配置}

发送天线数为8，接收天线数为32，迭代次数为2，Adder Tree流水级数为3，QAM阶数为16，时钟周期为\SI{10}{\nano\second}。

\subsection{预测公式}

\begin{align}
    L_{\mathrm{pred}} &= (2+1)(8+3)+4\times2-1=40\ \mathrm{cycles},\\
    I_{\mathrm{pred}} &= N_T=8\ \mathrm{cycles},\\
    B_{\mathrm{valid}} &= 8\log_2(16)=32\ \mathrm{bits},\\
    R_{\mathrm{pred}} &= \frac{32}{10\times8}=0.4\ \mathrm{Gbps},\\
    A_{\mathrm{pred}} &= F_{\mathrm{MIMO\_area}}(Config_{case1}),\\
    C_{\mathrm{pred}} &= \frac{A_{\mathrm{pred}}}{A_{\mathrm{GE}}}\times40.
\end{align}

\subsection{Validate过程}

Validate运行MIMO行为参考模型和RTL仿真，按照测试帧、发送天线和调制符号匹配完整检测向量；从VCD测量首个检测向量延迟和相邻向量输出间隔，并根据实测间隔计算参考Gbps。

\section{BP Polar Decoder模块Case 1}

\subsection{配置}

硬件架构为TypeI，译码算法为MS，码长为1024，并行度为256，数据位宽为5，码率为0.5，$E_b/N_0$为\SI{10}{\decibel}，时钟周期为\SI{20}{\nano\second}。

\subsection{预测公式}

\begin{align}
    G &= \frac{1024}{256}=4,\\
    C_{\mathrm{iter}} &= 2G\log_2(1024)-1=79\ \mathrm{cycles},\\
    C_{\mathrm{decision}} &= G=4\ \mathrm{cycles},\\
    L_{\mathrm{pred}} &= 79\operatorname{round}(I_{\mathrm{pred}})+4,\\
    B_{\mathrm{valid}} &= \operatorname{round}(1024\times0.5)=512\ \mathrm{bits},\\
    R_{\mathrm{pred}} &= \frac{512}{20L_{\mathrm{pred}}}\ \mathrm{Gbps},\\
    C_{\mathrm{pred}} &= \frac{A_{\mathrm{pred}}}{A_{\mathrm{GE}}}L_{\mathrm{pred}}.
\end{align}

\subsection{Validate过程}

Validate运行C/C++ BP译码程序，取得译码比特和实际迭代次数；运行RTL仿真并逐比特比较Hard Decision输出，同时检查冻结位、迭代结束状态和Early Stop结果；最后从波形取得实际译码延迟并计算参考bit吞吐率。

\chapter{结果记录和验收}

\section{单Case结果表}

\begin{table}[H]
    \centering
    \scriptsize
    \begin{adjustbox}{max width=\textwidth}
    \begin{tabular}{lcccccccc}
        \toprule
        \metricheader{指标} & \metricheader{单位} & \metricheader{预测值} & \metricheader{参考值} & \metricheader{误差} & \metricheader{预测时间} & \metricheader{参考时间} & \metricheader{提升倍数} & \metricheader{结论} \\
        \midrule
        面积 & $\mu\mathrm{m}^{2}$ & 待录入 & 待录入 & 待录入 & 待录入 & 待录入 & 待录入 & 待判定 \\
        延迟 & cycle & 待录入 & 待录入 & 待录入 & 待录入 & 待录入 & 待录入 & 待判定 \\
        吞吐率 & Gbps & 待录入 & 待录入 & 待录入 & 待录入 & 由RTL间隔计算 & -- & 待判定 \\
        硬件计算复杂度 & $\mathrm{GE\cdot cycle}$ & 待录入 & 待录入 & 待录入 & 待录入 & 由面积和延迟得到 & -- & 待判定 \\
        \bottomrule
    \end{tabular}
    \end{adjustbox}
    \caption{单Case指标记录模板}
\end{table}

\section{模块汇总表}

\begin{table}[H]
    \centering
    \scriptsize
    \begin{adjustbox}{max width=\textwidth}
    \begin{tabular}{lcccccccc}
        \toprule
        \metricheader{模块} & \metricheader{功能匹配} & \metricheader{自动评估时间} & \metricheader{整体提升倍数} & \metricheader{面积误差} & \metricheader{延迟误差} & \metricheader{吞吐率误差} & \metricheader{复杂度误差} & \metricheader{结论} \\
        \midrule
        Add & 待录入 & 待录入 & 待录入 & 待录入 & 待录入 & 待录入 & 待录入 & 待判定 \\
        Mul & 待录入 & 待录入 & 待录入 & 待录入 & 待录入 & 待录入 & 待录入 & 待判定 \\
        5G PUSCH CE & 待录入 & 待录入 & 待录入 & 待录入 & 待录入 & 待录入 & 待录入 & 待判定 \\
        INSA-MMSE MIMO Detector & 待录入 & 待录入 & 待录入 & 待录入 & 待录入 & 待录入 & 待录入 & 待判定 \\
        BP Polar Decoder & 待录入 & 待录入 & 待录入 & 待录入 & 待录入 & 待录入 & 待录入 & 待判定 \\
        \bottomrule
    \end{tabular}
    \end{adjustbox}
    \caption{已有模块指标汇总模板}
\end{table}

\section{验收规则}

\begin{enumerate}
    \item RTL输出必须与C/C++黄金输出一一对应；功能验证不通过时不判定指标误差。
    \item 自动评估总时间不超过\SI{10}{\second}。
    \item 相对于RTL验证和综合流程的整体速度提升不低于10倍。
    \item 基带电路自动评估误差不超过25\%。
    \item 硬件计算复杂度评估误差不超过15\%。
\end{enumerate}

\begin{thebibliography}{9}

\bibitem{nguyen2025mlkem}
T.-H. Nguyen, T.-K. Dang, D.-T. Dam, K.-D. Nguyen, P.-P. Duong,
C.-K. Pham, and T.-T. Hoang,
``An Area-Time Efficient Hardware Architecture for ML-KEM Post-Quantum
Cryptography Standard,'' \emph{IEEE Access}, vol. 13, pp. 103834--103847,
2025, doi: \href{https://doi.org/10.1109/ACCESS.2025.3579379}{10.1109/ACCESS.2025.3579379}.

\bibitem{curlin2025sbox}
P. S. Curlin, J. Heiges, C. Chan, and T. S. Lehman,
``A Survey of Hardware-Based AES SBoxes: Area, Performance, and Security,''
\emph{ACM Computing Surveys}, vol. 57, no. 9, pp. 1--37, 2025,
doi: \href{https://doi.org/10.1145/3724114}{10.1145/3724114}.

\end{thebibliography}

\end{document}
